% Part: model-theory
% Chapter: models-of-arithmetic
% Section: non-standard-models

\documentclass[../../../include/open-logic-section]{subfiles}

\begin{document}

% SEGMENT OLP-0195-S01

\olfileid{mod}{mar}{mdq}
\section{$\Th{Q}$ இன் மாதிரிகள்}

\begin{explain}
~$\Struct{N}$ மெய்யாக்கும் அதே !!{sentence}s ஐ மெய்யாக்கும், அதாவது
~$\Th{TA}$ இன் மாதிரியாக உள்ள, நிலையற்ற !!{structure}s உள்ளன என்பதை
அறிவோம். $\Sat{N}{\Th{Q}}$ என்பதால்,~$\Th{TA}$ இன் எந்த மாதிரியும்
~$\Th{Q}$ இன் மாதிரியும் ஆகும். $\Th{Q}$ ஆனது~$\Th{TA}$ ஐவிட
மிகவும் வலுக்குறைந்தது; எடுத்துக்காட்டாக,
$\Th{Q} \Proves/ \lforall[x][\lforall[y][\eq[(x +
      y)][(y+x)]]]$. வலுக்குறைந்த கோட்பாடுகளை நிறைவுறுத்துவது எளிது:
அவற்றுக்கு அதிக மாதிரிகள் உண்டு. எடுத்துக்காட்டாக,
$\lforall[x][\lforall[y][\eq[(x + y)][(y+x)]]]$ ஐப் பொய்யாக்கும்
மாதிரிகள் $\Th{Q}$ க்கு உள்ளன; ஆனால் அவை~$\Th{TA}$ இன் மாதிரிகளாகவோ,
அல்லது அதற்காகச் சொன்னால் $\Th{PA}$ இன் மாதிரிகளாகவோகூட இருக்க
முடியாது. $\Th{Q}$ இன் மாதிரிகள் ஒப்பளவில் எளிமையானவையும் ஆகும்:
அவற்றை வெளிப்படையாகக் குறிப்பிடலாம்.
\end{explain}

\begin{ex}
\ollabel{ex:model-K-of-Q}
ஆட்களம் $\Domain{K} = \Nat
\cup \{a\}$ உடைய !!{structure}~$\Struct{K}$ ஐயும் பின்வரும்
பொருள்கோள்களையும் கருதுக:
\begin{align*}
  \Assign{\Obj{0}}{K} & = 0\\
  \Assign{\prime}{K}(x) & =
  \begin{cases}
    x+1 & \text{if $x\in \Nat$}\\
    a & \text{if $x = a$}
  \end{cases}\\
  \Assign{+}{K}(x, y) & =
  \begin{cases}
    x+y & \text{if $x$, $y \in\Nat$}\\
    a & \text{otherwise}
  \end{cases}\\
  \Assign{\times}{K}(x, y) & =
  \begin{cases}
    xy & \text{if $x$, $y \in\Nat$}\\
    0 & \text{if $x = 0$ or $y = 0$}\\
    a & \text{otherwise}\\
  \end{cases}\\
  \Assign{<}{K} & =
  \Setabs{\tuple{x,y}}{x, y \in \Nat \text{ and } x<y} \cup
  \Setabs{\tuple{x,a}}{x \in \Domain{K}}
\end{align*}
$\Sat{K}{\Th{Q}}$ என்று காட்ட,~$\Th{Q}$ இன் எல்லா அடிகோள்களும்
~$\Struct{K}$ இல் மெய் என்று சரிபார்க்க வேண்டும். வசதிக்காக,
~$\Struct{K}$ இலுள்ள $x$ இன் ``அடுத்தெண்ணான''
$\Assign{\prime}{K}(x)$ க்கு $x^\nssucc$ என்றும்,~$\Struct{K}$
இலுள்ள $x$ மற்றும் $y$ இன் ``கூட்டுத்தொகையான''
$\Assign{+}{K}(x, y)$ க்கு $x \nsplus y$ என்றும்,~$\Struct{K}$
இலுள்ள $x$ மற்றும்~$y$ இன் ``பெருக்கற்பலனான''
$\Assign{\times}{K}(x, y)$ க்கு $x \nstimes y$ என்றும்,
$\tuple{x,y} \in
\Assign{<}{K}$ என்பதற்கு $x \nsless y$ என்றும் எழுதுவோம்.
இச்சுருக்கங்களால்~$\Struct{K}$ இலுள்ள செயல்களை இன்னும் தெளிவாக
பின்வருமாறு தரலாம்:
\[
\begin{array}{c|c}
  x & x^\nssucc \\
  \hline
  n & n+1 \\
  a & a
\end{array}
\qquad
\begin{array}{c|ccc}
  x \nsplus y & 0 & m & a \\
  \hline
  0 & 0 & m & a \\
  n & n & n+m & a \\
  a & a & a & a \\
\end{array}
\qquad
\begin{array}{c|ccc}
  x \nstimes y & 0 & m & a \\
  \hline
  0 & 0 & 0 & 0 \\
  n & 0 & nm & a \\
  a & 0 & a & a \\
\end{array}
\]
$n$, $m \in \Nat$ க்கு $n \nsless m$ எனில், எனில் மட்டுமே, $n<m$;
மேலும் எல்லா~$x \in \Domain{K}$ க்கும் $x \nsless
a$.

$\nssucc$ !!{injective} என்பதால்,
$\Sat{K}{\lforall[x][\lforall[y][(\eq[x'][y'] \lif \eq[x][y])]]}$.
~$\Struct{K}$ இல் $0$ ஓர் $\nssucc$-அடுத்தெண் அல்ல என்பதால்,
$\Sat{K}{\lforall[x][\eq/[\Obj
0][x']]}$. ஒவ்வொரு $n>0$ க்கும் $n = (n-1)^\nssucc$ என்றும்
$a = a^\nssucc$ என்றும் இருப்பதால்,
$\Sat{K}{\lforall[x][(\eq[x][\Obj 0] \lor \lexists[y][\eq[x][y']])]}$.

$n \nsplus 0 = n+0
= n$ மற்றும்~$\nsplus$ இன் வரையறைப்படி $a\nsplus 0 = a$ என்பதால்,
$\Sat{K}{\lforall[x][\eq[(x + \Obj 0)][x]]}$.
$\Sat{K}{\lforall[x][\lforall[y][\eq[(x + y')][(x + y)']]]}$ சற்றுச்
சிக்கலானது. $n$, $m$ ஆகிய இரண்டும் நிலையானவை எனில்:
\begin{align*}
(n \nsplus m^\nssucc) & = (n+(m+1)) = (n+m)+1 = (n \nsplus m)^\nssucc
\intertext{since $\nsplus$ and $^\nssucc$ agree with $+$ and $\prime$ on
  standard numbers.  Now suppose $x \in \Domain{K}$. Then}
(x \nsplus a^\nssucc) & = (x \nsplus a) = a = a^\nssucc = (x \nsplus a)^\nssucc
\intertext{The remaining case is if $y \in \Domain{K}$ but $x =
  a$. Here we also have to distinguish cases according to whether $y =
  n$ is standard or $y = b$:}
(a \nsplus n^\nssucc) & = (a \nsplus (n+1)) = a = a^\nssucc = (a \nsplus n)^\nssucc\\
(a \nsplus a^\nssucc) & = (a \nsplus a) = a = a^\nssucc = (a \nsplus a)^\nssucc
\end{align*}
இது தேவையைவிடச் சற்று விரிவானதே. எடுத்துக்காட்டாக, $z$ எதுவாக
இருந்தாலும் $a \nsplus z = a$ என்பதால், $a \nsplus
a^\nssucc = a$ என்று உடனே முடிவுகொள்ளலாம். மீதமுள்ள அடிகோள்களையும்
இதே வழியில் சரிபார்க்கலாம்.

எனவே $\Struct{K}$ ஆனது~$\Th{Q}$ இன் ஒரு மாதிரி. அதன்
``கூட்டல்''~$\nsplus$ பரிமாற்றுத்தன்மையுடையதும் ஆகும். ஆனால்
~$\Struct{N}$ இல் மெய்யாகவும்~$\Struct{K}$ இல் பொய்யாகவும் உள்ள
பிற வாக்கியங்களும், மறுதலையாக உள்ளவையும் உண்டு. எடுத்துக்காட்டாக,
$a \nsless a$; எனவே $\Sat{K}{\lexists[x][x < x]}$ மற்றும்
$\Sat/{K}{\lforall[x][\lnot x<x]}$. இதனால்
$\Th{Q} \Proves/
\lforall[x][\lnot x < x]$ என்று தெரிகிறது.
\end{ex}

\begin{prob}
\olref[mod][mar][mdq]{ex:model-K-of-Q} இலுள்ள $\Struct{K}$ ஆனது
~$\Th{Q}$ இன் மீதமுள்ள அடிகோள்களை நிறைவுறுத்துகிறது என்று நிரூபிக்க:
\begin{align*}
  & \lforall[x][\eq[(x \times \Obj 0)][\Obj 0]] \tag{$!Q_6$}\\
  & \lforall[x][\lforall[y][\eq[(x \times y')][((x \times y) + x)]]] \tag{$!Q_7$}\\
  & \lforall[x][\lforall[y][(x < y \liff \lexists[z][\eq[(z' + x)][y]])]] \tag{$!Q_8$}
\end{align*}
~$\Struct{N}$ இல் மெய்யாகவும்~$\Struct{K}$ இல் பொய்யாகவும் உள்ள,
~$\prime$ ஐ மட்டும் கொண்ட !!a{sentence} ஒன்றைக் கண்டறிக.
\end{prob}

\begin{ex}
\ollabel{ex:model-L-of-Q} ஆட்களம் $\Domain{L} = \Nat \cup \{a, b\}$
உடைய !!{structure}~$\Struct{L}$ ஐக் கருதுக. பொருள்கோள்கள்
$\Assign{\prime}{L} = \nssucc$, $\Assign{+}{L} = \nsplus$ பின்வருமாறு:
\[
\begin{array}{c|c}
  x & x^\nssucc \\
  \hline
  n & n+1 \\
  a & a\\
  b & b
\end{array}
\qquad
\begin{array}{c|ccc}
  x \nsplus y & m & a & b\\
  \hline
  n & n+m & b & a\\
  a & a & b & a\\
  b & b & b & a
\end{array}
\]
$\nssucc$ !!{injective} ஆகவும், $0$ அதன் வீச்சில் இல்லாமலும்,
~$0$ தவிர ஒவ்வொரு $x \in \Domain{L}$ உம் அதன் வீச்சில் இருப்பதாலும்,
அடிகோள்கள் $!Q_1$--$!Q_3$ ஆனவை~$\Struct{L}$ இல் மெய். எந்த
$x$ க்கும் $x \nsplus 0 = x$; எனவே $!Q_4$ உம் மெய். $!Q_5$ க்கு,
$x \nsplus y^\nssucc$ மற்றும் $(x \nsplus
y)^\nssucc$ ஐக் கருதுக. $x$ மற்றும் $y$ இரண்டும் நிலையானவை எனில்,
$\nssucc$ மற்றும் $\nsplus$ ஆனவை $\prime$ மற்றும் $+$ உடன்
உடன்படுவதால் அவை சமம். $x$ நிலையற்றதாகவும் $y$ நிலையானதாகவும்
இருந்தால், $x \nsplus y^\nssucc = x =
x^\nssucc = (x \nsplus y)^\nssucc$. $x$ மற்றும் $y$ இரண்டும்
நிலையற்றவை எனில், நான்கு நிலைகள் உள்ளன:
\begin{align*}
&  a \nsplus a^\nssucc  = b = b^\nssucc = (a \nsplus a)^\nssucc\\
&  b \nsplus b^\nssucc  = a = a^\nssucc = (b \nsplus b)^\nssucc\\
&  b \nsplus a^\nssucc  = b = b^\nssucc = (b \nsplus y)^\nssucc\\
&  a \nsplus b^\nssucc  = a = a^\nssucc = (a \nsplus b)^\nssucc\\
\intertext{If $x$ is standard, but $y$ is non-standard, we have}
&  n \nsplus a^\nssucc  = n \nsplus a = b = b^\nssucc = (n \nsplus a)^\nssucc\\
&  n \nsplus b^\nssucc  = n \nsplus b = a = a^\nssucc = (n \nsplus b)^\nssucc
\end{align*}
எனவே $\Sat{L}{!Q_5}$. ஆனால் $a \nsplus 0 \neq 0 \nsplus a$;
எனவே $\Sat/{L}{\lforall[x][\lforall[y][\eq[(x+y)][(y+x)]]]}$.
\end{ex}

\begin{prob}
\olref[mod][mar][mdq]{ex:model-L-of-Q} இலுள்ள $\Struct{L}$ உடன்,
~$\times$ மற்றும் $<$ ஐப் பொருள்கொள்ளும் $\nstimes$ மற்றும்
$\nsless$ ஐச் சேர்த்து விரிவாக்குக. உங்கள் கட்டமைப்பு~$\Th{Q}$ இன்
மீதமுள்ள அடிகோள்களை நிறைவுறுத்துகிறது என்று காட்டுக:
\begin{align*}
& \lforall[x][\eq[(x \times \Obj 0)][\Obj 0]] \tag{$!Q_6$}\\ &
  \lforall[x][\lforall[y][\eq[(x \times y')][((x \times y) + x)]]]
  \tag{$!Q_7$}\\ & \lforall[x][\lforall[y][(x < y \liff \lexists[z][\eq[(z'+x)][y]])]] \tag{$!Q_8$}
\end{align*}
\end{prob}

\begin{prob}
\olref[mod][mar][mdq]{ex:model-L-of-Q} இலுள்ள $\Struct{L}$ இல்,
$a^\nssucc
= a$ மற்றும் $b^\nssucc = b$. $a^\nssucc = b$ மற்றும்
$b^\nssucc = a$ ஆகும்~$\Th{Q}$ இன் மாதிரி ஒன்று உள்ளதா?
\end{prob}

\begin{explain}
நிலையற்ற !!{element}s, நிலையான உறுப்புகளுக்கு ``அப்பால்'' வாழும்
~$\Th{Q}$ இன் மாதிரிகளை வெளிப்படையாகக் கட்டமைத்துள்ளோம். உண்மையில்,
அடிகோள்கள் இதைத் தேவைப்படுத்துகின்றன. நிலையற்ற !!{element}~$x$
ஆனது ${} \nsless 0$ ஆக முடியாது; ஏனெனில்
$\Th{Q} \Proves
\lforall[x][\lnot x<0]$ (\olref[inc][req][min]{lem:less-zero} ஐப்
பார்க்க). மேலும், ஒவ்வொரு $n$ க்கும்,
$\Th{Q} \Proves \lforall[x][(x < \num{n}' \lif
(\eq[x][\num{0}] \lor \eq[x][\num{1}] \lor \dots \lor
\eq[x][\num{n}]))]$ (\olref[inc][req][min]{lem:less-nsucc}); எனவே
எந்த~$n>0$ க்கும் $a \nsless n$ ஆக முடியாது.
\end{explain}

\end{document}
